\documentclass[10pt,a4paper]{amsart}
\usepackage[margin=19mm]{geometry}
\usepackage{amsmath,amssymb,mathtools}
\usepackage[scr=rsfs,cal=euler]{mathalfa}
\usepackage{xfrac}
\usepackage[unicode,hidelinks,bookmarks=false]{hyperref}
\newcommand{\ie}{i.e.\ }
\newcommand{\cf}{cf.\ }
\newcommand{\resp}{respectively\ }
\newcommand{\Subsection}{\subsection}
\providecommand{\nobreakdash}{\nobreak\hbox{-}}
\setlength{\emergencystretch}{2em}
\multlinegap0pt
\usepackage{amssymb}
\newtheorem{lemma}{Lemma}
\newtheorem{theorem}{Theorem}
\newtheorem{proposition}{Proposition}
\usepackage{cleveref}
\crefname{lemma}{Lemma}{Lemmas}
\crefname{theorem}{Theorem}{Theorems}
\crefname{proposition}{Proposition}{Propositions}
\DeclareMathOperator{\Span}{\mathrm{span}}
\DeclareMathOperator{\aff}{aff}
\DeclareMathOperator{\cone}{cone}
\DeclareMathOperator{\convcone}{convcone}
\DeclareMathOperator{\lin}{lin}
\title{Convexity from order: a theory of faces and generalized separation}
\author{Matthew S. Scott}
\date{Source: arXiv:2609.05775v1 (CC BY 4.0)}
\begin{document}
\maketitle

\noindent\emph{Source and licence.} Convexity from order: a theory of faces and generalized separation. Version arXiv:2609.05775v1 (CC BY 4.0), distributed by arXiv under the Creative Commons Attribution 4.0 International licence, http://creativecommons.org/licenses/by/4.0/, as recorded in the arXiv licence metadata for this exact version and shown on the abstract page https://arxiv.org/abs/2609.05775v1. Full licence notice: \texttt{licenses/arxiv-2609.05775-CC-BY-4.0.txt}.\par\smallskip

\noindent\textit{Editorial scope.} Counted unit: the article's theorem loc:affine\_cross:lineality.statement (source0.tex lines 1415--1428). The article's complete original proof of it is delivered with it (source0.tex lines 1497--1543, one \texttt{proof} environment of 47 lines, which the article places away from the statement). No mathematics is written, repaired or re-derived: the statement and the proof are the article's own lines, in the article's own notation, and no retained line is substituted or re-flowed.  Retained uncounted as the source-local context the proof needs: lines 688--694; lines 1048--1064; lines 1168--1178; lines 1247--1256; lines 1274--1280; lines 1319--1329; lines 1475--1481.  Nothing was omitted from any retained span, so the package carries no omission record.  No retained line contains a citation, so the wrapper carries no bibliography and no bibliography entry.  No retained line contains a figure environment, an image-inclusion command or drawing code, so no image is delivered and the package declares no asset dependency.  Editorial formatting only, in addition: an AMS \texttt{amsart} preamble that restates the article's own declarations and macros used by the retained lines, namely the environments \texttt{lemma}, \texttt{theorem}, \texttt{proposition} on separate counters, together with the standard packages those lines need.  The retained environments print in this document as Lemma 1, Theorem 1, Proposition 1 in order of appearance, whereas the article prints the same items with the numbering of its own full document.  The complete native source of the article as distributed in the arXiv e-print for this exact version is bundled unchanged as \texttt{sources/209463/source0.tex}.
\begin{lemma}[Span dehomogenizes to the affine hull]
\label{loc:span_dehomogenizes_to_the_affine_hull.statement}
Let $S \subseteq X$ be a convex set and $H := \cone \circ \gamma_1(S)$ its homogenization cone. Then
\begin{equation*}
\gamma_1^{-1}(\Span H) = \aff(S).
\end{equation*}
\end{lemma}

\begin{theorem}[Characterization of the separating flats]
\label{loc:characterization_of_the_separating_flats.statement}
Given convex sets $C, D \subseteq X$, an affine flat $L \subseteq X$ is a separating flat when,
equivalently,
\begin{enumerate}
\item $L = \gamma_1^{-1}(U)$ for $U$ a separating subspace of $K = \cone \circ \gamma_1(C)$ and $G = \cone \circ \gamma_1(D)$.
\item $L = \ker f$ for some affine oriented $f:X \to (Y, Y_+)$ separating $C$ from $D$.
\end{enumerate}

And if $C \cap D \neq \varnothing$, the following statements are also equivalent.

\begin{enumerate}
\setcounter{enumi}{2}
\item $L$ is supporting for both $C$ and $D$, and $\forall c  \in C, d  \in D, c-d \in L-L \implies c,d \in L$.
\item $L =  \tilde{L} + C \cap D$ for $\tilde{L}$ a supporting subspace of $\cone(C-D)$.
\end{enumerate}
\end{theorem}

\begin{lemma}[Characterization of the lineality]
\label{loc:characterization_of_the_lineality.statement}
Given a convex cone $K \subseteq W$, the lineality $\lin(K)$ is, equivalently,
\begin{enumerate}
\item $K \cap -K$
\item $\{x \sim_{K} 0 \mid x \in K\}.$
\item The largest subspace contained in $K$; it contains every subspace contained in $K$. 
\item The minimum supporting subspace of $K$.
\item The minimum face of $K$.
\end{enumerate}
\end{lemma}

\begin{theorem}[Characterization of the cross-lineality]
\label{loc:cross_lineality.statement}
For two convex cones $K, G \subseteq W$, the \emph{cross-lineality} of the pair $(K, G)$ is equivalently defined as
\begin{enumerate}
\item $(K-G) \cap (G -K)=\lin(K-G)=\lin(G-K).$
\item The minimum separating subspace of $K$ and $G$.
\item $H_K(G)+ H_G(K)$
\item $F_K(G) - F_G(K)$
\end{enumerate}
\end{theorem}

\begin{proposition}[Generated face spans the lineality]
\label{loc:generated_face_spans_the_lineality.statement}
For convex cones $K, G \subseteq W$,
\begin{equation*}
H_K(G) = \Span(F_K(G)).
\end{equation*}
\end{proposition}

\begin{theorem}[Characterization of the generated face of a convex set]
\label{loc:characterization_of_the_generated_face_of_a_convex_set.statement}
For convex sets $C, D \subseteq X$ with homogenization cones $K, G \subseteq X \times \mathbb{R}$, the generated face 
$F_C(D)$ is equivalently:
\begin{enumerate}
\item The smallest face of $C$ that contains $C \cap D$.
\item $\gamma_1^{-1}(F_K(G))$.
\item $\min_{\cone(C-D)}C$ when $C \cap D \neq \varnothing$, otherwise it is $\varnothing$.
\item $\min_{\cone(C- (C \cap D))}C$ when $C \cap D \neq \varnothing$, otherwise it is $\varnothing$.
\end{enumerate}
\end{theorem}

\begin{proposition}[Generated faces are formed by symmetric differences]
\label{loc:generated_faces_are_formed_by_symmetric_differences.statement}
For convex sets $C, D \subseteq X$, the points in $F_C(D)$ are precisely those that are part of a symmetric difference:
\begin{equation*}
F_C(D) = \{c \in C : c-d = -\lambda(c'-d') \text{ for some } \lambda > 0,\ c' \in C,\ d, d' \in D\}.
\end{equation*}
\end{proposition}

\begin{theorem}[Characterizations of the joint supporting subspace]
\label{loc:affine_cross:lineality.statement}
Let $C, D \subseteq X$ be convex sets,
with their homogenization cones $K = \cone(C \times \{ 1 \})$ and $G = \cone(D \times \{ 1 \})$.
The \emph{joint supporting subspace} $T_a(C, D)$ is equivalently:
\begin{enumerate}
\item $\gamma_1^{-1}(\lin(K-G))$
\item $\lin\cone(C-D) + C \cap D$
\item The minimum separating flat of $C$ and $D$.
\item $\lin \cone(C - (C \cap D)) + \lin \cone(D - (C \cap D)) + C \cap D$.
\item $\aff(F_C(D) \cup F_D(C))$
\item $\aff\{ c,c',d,d'\mid c, c'  \in C , d,d'  \in D \text{ s.t. } c-d = \lambda(d'-c') \text{ for some } \lambda > 0\}$
\end{enumerate}
\end{theorem}

\begin{proof}[\hypertarget{loc:affine_cross:lineality.proof}{}Proof of \Cref{loc:affine_cross:lineality.statement}]

If $C \cap D =  \emptyset$, then all points are the same set because they are
all the empty set. Next we show equality between all points in the case that $C \cap D \neq \varnothing$.

$1 = 3$:
By \Cref{loc:characterization_of_the_separating_flats.statement} point 1,
the separating flats of $(C,D)$ are exactly the preimages $\gamma_1^{-1}(U)$ of the separating subspaces $U$ of $(K,G)$. By \Cref{loc:cross_lineality.statement} point 2, $\lin(K-G)$ is the minimum separating subspace. Since $\gamma_1^{-1}$ preserves inclusion, it sends this minimum to the minimum separating flat, which is point 1. We note this argument holds even when $C \cap D = \varnothing$.

$3 = 2$:
By \Cref{loc:characterization_of_the_separating_flats.statement} point 4, the map $\tilde{L} \mapsto \tilde{L} + C \cap D$ sends the supporting subspaces of $\cone(C-D)$ onto the separating flats of $(C,D)$, and it preserves inclusion. By \Cref{loc:characterization_of_the_lineality.statement} point 4, $\lin\cone(C-D)$ is the minimum supporting subspace, so its image $\lin\cone(C-D) + C \cap D$ is the minimum separating flat.

$1 = 5$:
We compute from point 1:
\begin{align*}
\gamma_1^{-1}(\lin(K-G))
&= \gamma_1^{-1}\Span\bigl(F_K(G) \cup F_G(K)\bigr)\\
&= \gamma_1^{-1}\Span\bigl(\convcone(F_K(G) \cup F_G(K))\bigr)\\
&= \aff \gamma_1^{-1}\bigl(\convcone(F_K(G) \cup F_G(K))\bigr)\\
&= \aff\bigl(\gamma_1^{-1}(F_K(G)) \cup \gamma_1^{-1}(F_G(K))\bigr)\\
&= \aff\bigl(F_C(D) \cup F_D(C)\bigr),
\end{align*}
which is point 5. The first equality is \Cref{loc:cross_lineality.statement} point 3 with the identification of $H_K(G) = \Span(F_K(G))$ by \Cref{loc:generated_face_spans_the_lineality.statement}, with the fact that 
$\Span(F_K(G)) + \Span(F_G(K)) =  \Span(F_K(G) \cup F_G(K))$. The third is \Cref{loc:span_dehomogenizes_to_the_affine_hull.statement}, applied to the homogenization cone $\convcone(F_K(G) \cup F_G(K)) \subseteq (X \times \mathbb{R}_{++}) \cup \{0\}$.
The fifth is \Cref{loc:characterization_of_the_generated_face_of_a_convex_set.statement} point 2.

$5 = 4$:
First consider that
\begin{align*}
\aff(F_C(D)) &= \aff\bigl(F_C(D) \cup (C \cap D)\bigr) \\
&= \aff\bigl(F_C(C \cap D) \cup F_{C \cap D}(C)\bigr) \\
&= \lin\cone(C - (C \cap D)) + C \cap D =: L_C + C \cap D,
\end{align*}
where we defined $L_C := \lin \cone(C-(C \cap D))$ in the last line. The first equality holds because $C \cap D \subseteq F_C(D)$ (follows from \Cref{loc:characterization_of_the_generated_face_of_a_convex_set.statement} point 4). The second equality uses the fact that $F_C(C \cap D) = F_C(D)$ by equivalence of points 3 and 4 of \Cref{loc:characterization_of_the_generated_face_of_a_convex_set.statement}, and also that $F_{C \cap D}(C) = C \cap D$ because $C \cap D \subseteq C$. The third equality uses the equivalence between
points 5 and 2 shown above, applied to the convex sets $C$ and $C \cap D$.

Similarly, we have $\aff(F_D(C)) = L_D + C \cap D$. The subspaces $L_C$ and $L_D$ both contain $C \cap D - C \cap D$, therefore the term $(+ C \cap D)$ is a common affine offset applied to both. Therefore their affine hull is the flat of direction $L_C + L_D$ through $C \cap D$, i.e.,
\begin{equation*}
\aff(F_C(D) \cup F_D(C)) = L_C + L_D + C \cap D,
\end{equation*}
which is point 4.

$6 = 5$:
By \Cref{loc:generated_faces_are_formed_by_symmetric_differences.statement},
the set of points forming symmetric differences is
precisely $F_C(D) \cup F_D(C)$.
\end{proof}

\end{document}
